\subsection{Structure of the Thesis}
This thesis is divided into ten chapters, each building upon the previous to explore the use of AI-based methods—specifically LSTM architectures—for anomaly detection in Electric Vehicle Supply Equipment (EVSE).\\
Chapter~1 introduces the motivation and relevance of the topic, outlines the role of anomaly detection in the context of EV infrastructure, and formulates the central research questions. Additionally, it provides an overview of prior AI-based approaches and concludes with this structural outline. Chapter~2 summarizes the key theoretical foundations necessary for the thesis, including the basics of anomaly detection, a short introduction to tree-based models such as decision trees and random forests, and a focused explanation of the LSTM architecture.\\
Chapter~3 presents related work, covering research on security aspects and anomaly detection mechanisms specific to EVSE, as well as a broader review of AI-based detection methods applied in similar domains. Chapter~4 describes the dataset creation and acquisition process, detailing both public datasets and custom setups involving Prometheus monitoring tools. Chapter~5 follows with an exploratory data analysis (EDA), intended to uncover key patterns and differences between the public and self-generated datasets.\\
In Chapter~6, the process of data preprocessing and feature selection is discussed, including the identification of relevant features and the use of feature importance methods such as decision trees. Chapter~7 focuses on the LSTM model itself, describing its architecture, the training process, and the chosen strategies for hyperparameter tuning. \\
The results of the approach are then evaluated in Chapter~8, which outlines the experimental setup, evaluation metrics, and key results of the applied anomaly detection system. Chapter~9 provides a detailed discussion of the findings in relation to the research questions, highlighting limitations and potential practical applications. Finally, Chapter~10 concludes the thesis with a summary of the work and a reflection on future directions.\\
This structure ensures a logical and methodical development of the topic, from theoretical foundations to implementation, evaluation, and critical interpretation.
